% Chapter 8 — Validación
\chapter{Validación}

\section{Estrategia}

Comparación de la línea de costa detectada automáticamente contra
\textbf{transectos GNSS independientes} del IEL, es decir, puntos
\textbf{no usados} en la calibración de homografía.

\section{Métricas}

\begin{description}
  \item[RMSE] Error cuadrático medio:
        \[
          \text{RMSE} = \sqrt{\frac{1}{N}\sum_{i=1}^N d_i^2}
        \]
        donde $d_i$ es la distancia mínima (metros) del punto GNSS $i$
        a la polilínea detectada.

  \item[MAE] Error medio absoluto:
        \[
          \text{MAE} = \frac{1}{N}\sum_{i=1}^N d_i
        \]
\end{description}

\section{Criterios de aceptación}

\begin{itemize}
  \item RMSE $< 1.5$ m en condiciones estándar (12:00h, sin interferencias).
  \item MAE $< 1.0$ m como objetivo secundario.
  \item Cobertura temporal: $\geq 90\%$ de imágenes a 12:00h procesadas
        sin error.
\end{itemize}

\section{Factores de error a monitorizar}

\begin{itemize}
  \item Calidad de imagen: niebla, reflejo solar, lluvia.
  \item Mareas: afectan la posición real de la línea de costa.
  \item Deriva temporal de la cámara (movimiento del soporte).
  \item Precisión de los GCPs GNSS de referencia.
\end{itemize}
